\section{Context Engineering}
Bien que les modèles récents acceptent des contextes de plus en plus larges (jusqu'à 128k tokens), injecter l'intégralité du code source d'un projet dans le prompt est une stratégie inefficace. Non seulement cela augmente drastiquement la latence et les coûts, mais cela dilue l'information pertinente dans du bruit ("Lost in the Middle phenomenon") \cite{Liu2023}. Il est donc primordial de gérer correctement le contexte fourni au LLM.

\subsection{Les instructions de base}
Le système de \textit{context engineering} s'appuie sur un prompt fondateur qui guide tout le raisonnement du LLM. Il définit la philosophie à suivre par le biais des consignes censées orienter le raisonnement du modèle. Par exemple, on indique au LLM de chercher en priorité dans les scripts et non dans les fichiers de données qui sont bien souvent moins pertinents pour répondre à une question sur une mécanique globale.

\begin{lstlisting}
Critical Mental Model
    When a user asks "where is X configured?" or "how do I modify Y?", the safest default is to look for the script that WRITES the relevant table or file.
    If a data file is only being READ, it is often downstream consumption; the business logic is usually upstream in the script that produced it.
\end{lstlisting}

Ce document définit également des connaissances essentielles sur la syntaxe Envision pour permettre au LLM de mieux identifier différentes sections du code. Le langage Envision étant un langage interne à Lokad, cette section est absolument essentielle. En effet, le LLM peut procéder par analogie avec d'autres langages plus populaires ou comprendre grâce au contexte mais il demeure qu'il n'a pas été entraîné dessus auparavant et n'en maîtrise donc pas la syntaxe.

\begin{lstlisting}
Variables and Tables
    Vectors vs. Scalars: * Vectors (columns): Identified by dot notation (e.g., Orders.Amount).
    Scalars (single values): Identified without a dot (e.g., exchangeRate).
    Dimensions: Tables typically define a primary key within brackets (e.g., Items[id]).
\end{lstlisting}

Le document \texttt{base\_instructions.txt} est unique et utilisé à chaque étape du \textit{planning} afin d'assurer un raisonnement cohérent entre chaque appel.

\subsection{Le Distillation Tool}
Une fois le retrieval effectué, nous obtenons des \texttt{RetrievalResults} qui contiennent des chunks entiers, dont le volume d'information apparait superflu par rapport à la question initiale.La mission du distillation tool est de condenser ces \texttt{RetrievalResults} en des faits concis, en rapport avec la question. Dans le cas où un fait est redondant entre les chunks, on n'en garde alors qu'un seul exemplaire.

Notons que nous avons recours à de la \textit{Lazy Distillation} : nous n'effectuons pas la phase de distillation immédiatement après la récupération d'informations par un outil, mais plutôt lors de la phase de planning qui suit. Cela permet d'éviter de condenser des \texttt{RetrievalResults} si le planner décide de rendre une réponse définitive à l'utilisateur.

\begin{lstlisting}
_DISTILLATION_SYSTEM_PROMPT = (
    "You are the Memory Manager of a complex RAG Agent. "
    "Your sole responsibility is to extract and preserve key facts "
    "so that the agent can answer its query without re-reading those documents.")
\end{lstlisting}

Cette distillation a lieu entre les différents appels au \textit{Planner}, avec accès aux  résultats bruts et à la tentative précédente du LLM principal afin de compresser au mieux l'information utile.

\subsection{Le Prior Evidence Tool}
Le \texttt{prior\_evidence\_tool} est un autre élément important du \textit{context engineering}. Il permet au \textit{Planner} de réaccéder à des informations précédentes via un système d'Id sans relancer une recherche par les \textit{tools}. 
Cela permet d'économiser des recherches inutiles, de combiner des informations issues d'appels différents et d'accéder à des faits issus de questions précédentes lorsque l'utilisateur est en mode conversationnel.

\subsection{Knowledge Bank / Accumulated Evidence / Execution History}
Nous utilisons trois "banques" de données qui permettent de conserver l'information et le contexte des opérations déjà réalisées. Chacune de ces banques possède une utilité différente, comme détaillé par le tableau suivant.

\renewcommand{\arraystretch}{1.5}
\begin{table}[h]
\centering
\label{tab:data_structures}
\begin{tabular}{|m{3cm}|m{3cm}|m{4cm}|m{4cm}|}
\hline
\rowcolor{gray!30} \textbf{Nom} & \textbf{Type de Donnée} & \textbf{Objectif} & \textbf{Source} \\ \hline
Knowledge Bank & List[Knowledge Element] & Faits distillés donnés en contexte au \textit{Planner} & Incrémenté par le \textit{Distillation Tool}\\ \hline
Accumulated Evidence & Dict[Id $\rightarrow$ Retrieval Result] & Banque de résultats bruts accessible via le \texttt{prior\_evidence\_tool} & Incrémenté par le \textit{Distillation Tool} \\ \hline
Execution History & List[ActionLog] & Résumé des appels précédents fourni au \textit{Planner} pour éviter la redondance. & Incrémenté directement par les \textit{tools}  \\ \hline
\end{tabular}
\caption{Synthèse des structures de données du projet}
\end{table}

Ici, il est important de noter que la \texttt{Knowledge Bank} ainsi qu'une version condensée de \texttt{l'Execution History} sont fournis dans le \textit{planner prompt} avec les mentions \textit{VERIFIED FACTS} et \textit{HISTORY}. En aucun cas le \textit{planner} ne reçoit les données brutes non condensées comme celles que contient \texttt{Accumulated Evidence}. Cette dernière banque fait office de base de données consultable au travers du \texttt{prior\_evidence\_tool} par le \textit{Planner} s'il en prend la décision.

\begin{figure}[htbp]
    \centering
\begin{tikzpicture}[
    node distance=1.2cm and 0.5cm,
    % Styles unifiés
    process/.style={
        rectangle, draw=blue!50, fill=blue!5, rounded corners, 
        align=left, text width=6.5cm, font=\small, drop shadow
    },
    step/.style={
        rectangle, draw=gray!80, fill=white, 
        text width=6.5cm, align=center, font=\bfseries\footnotesize % text width + align = résout l'erreur
    },
    arrow/.style={-{Stealth[scale=1.2]}, thick, gray}
]

% Nodes
\node[step] (start) {DÉBUT : Utilisateur pose une question};

\node[process, below=of start] (kickoff) {
    \textbf{DÉBUT}\\
    - Le \textit{Planner} lit \textit{base\_instructions} + \\ \textit{knowledge\_bank} (vide)\\
    - Sélection du premier outil (ex: \textit{rag\_tool})
};

\node[process, below=of kickoff] (rag) {
    \textbf{EXÉCUTION DU TOOL CHOISI}\\
    - Ex: Récupération de 5 documents via \textit{rag\_tool} \\
    - Stockage dans \textit{accumulated\_evidence}\\
    - Mise à jour de \textit{execution\_history}
};

\node[process, below=of rag] (distill) {
    \textbf{DISTILLATION}\\
    - Lecture des résultats bruts\\
    - Appel \textit{distill\_batch()} $\rightarrow$ 3 faits distillés\\
    - Mise à jour de \textit{knowledge\_bank}
};

\node[process, below=of distill] (continue) {
    \textbf{PLANNER PROMPT}\\
    - Lecture \textit{knowledge\_bank} + \textit{execution\_history}\\
    - Construction du prompt (Faits + Historique)\\          
};

\node[process, below=of continue] (decision) {
    \textbf{PLANNER DECISION} \\
    Différents choix possibles : \\
    -Soumettre une réponse définitive, \\
    -Combiner des faits ensemble \\
    -\textit{Retrieve} de nouveaux faits
};


\node[step, below=of decision] (end) {...};

% Arrows
\draw[arrow] (start) -- (kickoff);
\draw[arrow] (kickoff) -- (rag);
\draw[arrow] (rag) -- (distill);
\draw[arrow] (distill) -- (continue);
\draw[arrow] (continue) -- (decision);
\draw[arrow] (decision) -- (end);

\end{tikzpicture}
\caption{Exemple illustratif de la gestion du contexte (sans \textit{Lazy Distillation})}
\end{figure}

\subsection{Motivations derrière cette implémentation}
Les motivations derrière cette implémentation sont multiples :
\begin{itemize}
    \item \textbf{Efficacité en termes de tokens :} on réduit la taille du contexte en le limitant aux faits pertinents liés à la question. On évite des appels répétitifs grâce à l'historique des appels.
    \item \textbf{Traçabilité du raisonnement :} Ce système permet de suivre le cheminement du raisonnement via le système d'\texttt{ActionLogs} et d'\textit{Ids} liés aux faits. Cela permet à l'utilisateur de comprendre le cheminement et au développeur de debugger plus aisément.
    \item \textbf{Mode conversationnel avec mémoire :} Puisqu'on garde la trace des appels précédents, l'utilisateur peut exploiter ses anciennes questions lors d'une conversation avec le LLM.
\end{itemize}